\documentclass[../thesis.tex]{subfiles}
% \graphicspath{{\subfix{../images/}}}

\begin{document}

\chapter{Unsupervised domain adaptation}

\paragraph{Abstract.} This chapter explores how the group-instance model can bring improvements to the task of \textit{unsupervised domain adaptation}. This is the task of training a predictor function to perform well across multiple data environments with different relationships between covariate and response. My purpose in this chapter is to advance the conversation of the domain adaptation research community from the current evaluation paradigm of \textit{singleton prediction} (i.e., measuring the prediction performance on a single covariate) to \textit{contextual prediction} (i.e., measuring the prediction performance on a group of covariates from the same environment).

\section{Introduction}

The goal of unsupervised domain adaptation is to learn predictors that can be applied to various real-world data environments. The challenge is that, across environments, the relationship between covariates, $X$, and responses, $Y$, may differ. For example, consider the task of predicting the duration of a hike based on its distance. The same hiking trail (covariate) might take two hours to complete (response) for an experienced hiker (environment 1) but might take three hours to complete for a beginner (environment 2). The hope is that we can learn a predictor that performs well even on covariate-response pairs from a new environment (e.g., a third hiker about whose experience we know nothing).

The standard setup for domain adaptation involves training the predictor on a dataset comprising samples from a few different environments, and then evaluating the predictor on one sample coming from a new environment. Formally, let $\myvect{x}^i_{k} \in \mathbb{R}^{\mydim{x}}$ be the $k$-th covariate in environment $i$ of the training set, with $i \in \{1,\dots,\mycnt{n}\}$ and $k \in \{1,\dots,\mycnt{m}_i\}$, where $\mycnt{m}_i$ is the number of samples from environment $i$. Let $\myvect{y}^i_{k} \in \mathbb{R}^{\mydim{y}}$ be the $k$-th response in environment $i$. Let $\left(\myvect{x}', \myvect{y}' \right)$ be an covariate-response pair from the testing environment. The predictor uses the training set to learn the relationship between covariates and responses, and its prediction accuracy is then evaluated on the testing pair. I call this evaluation paradigm \textit{singleton prediction}. Actually, the predictor is typically evaluated independently on multiple testing pairs sampled from one or more new environments, and then the average prediction accuracy over all pairs is reported. However, since predictions are made independently for each testing pair, this is still a case of singleton prediction.

I am proposing an alternative evaluation setup where, instead of a single sample, the predictor can see a group of samples from a new testing environment, $\left\{\myvect{x}'_{k}, \myvect{y}_{k}\right\}_{k=1}^{\mycnt{m}'}$, and can make a joint prediction for the entire group. I call this paradigm \textit{contextual prediction}.

I claim that the new evaluation task of contextual prediction is more interesting than singleton prediction because (1) it allows for new solutions to the problem of domain adaptation (2) while still corresponding to a realistic evaluation scenario.

\section{Illustrative example: Diabetes risk}
\label{sec:example-diabetes}

\begin{figure}[t]
    \begin{center}
        \centerline{\includegraphics[width=0.5\columnwidth]{figures/prob_model.pdf}}
        \caption{Causal diagram for the variables involved in the diabetes risk model. A major cause of type-2 diabetes, $\myrv{y}$, is a high body-mass index, $\myrv{x}_{c}$, which itself is affected by whether the patient has a diet rich in salt, $W$. The frequency of thirst, $\myrv{x}_{s}$, is symptom of both diabetes, $\myrv{y}$, and a salt-rich diet, $\myrv{w}$. The environments in this case correspond to different groups of patients organised by diet.}
        \label{fig:diabetes}
    \end{center}
    \vskip -0.2in
\end{figure}

What follows is a concrete case study that exemplifies how the paradigm of contextual prediction is more interesting and relevant for domain adaptation. Consider the task of screening patients for type-2 diabetes risk. Let us imagine that we are developing a non-invasive test that predicts the risk of diabetes, $\myrv{y}$, based on body-mass index, $\myrv{x}_{c}$ (normalised to a real number between 0 and 1), and frequency of thirst (also a real number between 0 and 1), $\myrv{x}_{s}$. These variables could be measured through a simple questionnaire.

Even though both body-mass index and thirst frequency are predictive of diabetes, only body-mass index is a reliable predictor across environments \citep{ruigePerformanceNIDDMScreening1997}. This is because body-mass index is a cause of type-2 diabetes, whereas frequent thirst is a symptom. Causes make for more reliable predictors than symptoms because, when you intervene upon a cause, the symptom changes accordingly, but when you intervene upon an effect, the cause remains unchanged. In the case of domain adaptation, the different data environments reflect different interventions applied to the variables of the system \citep{petersCausalInferenceUsing2016}.

Suppose that the interventions in our diabetes example result from the amount of salt in the patient's diet, represented by the tertiary variable $\myrv{w}$. Salt-intake has a direct causal effect on both body-mass index, $\myrv{x}_{c}$, and on the frequency of thirst, $\myrv{x}_{s}$, but not on diabetes, $\myrv{y}$. The value of this variable is hidden to us (we do not measure it), but our patients are grouped according to diet (e.g., mediterranean, vegetarian, etc.) so we know that the average salt intake is different in each group.

Let us put numbers to our diabetes example in order to see why causes are consistent predictors. Consider the following causal mechanism for the variables involved in diabetes risk:
\begin{equation}\label{eq:diabetes-case-study}
\begin{split}
    \myrv{w} &\in \left[ 0,1,2 \right] \\
    \myrv{x}_{c} &= \mathrm{Beta} \Biggl( \alpha = 2 + \frac{1}{2}  \myrv{w}, ~ \beta = 2 - \frac{1}{2}  \myrv{w} \Biggr) \\
    \myrv{y} &= \mathrm{Bin} \Biggl( n = 1, ~ p = \myrv{x}_{c} \Biggr) \\
    \myrv{x}_{s} &= \mathrm{Beta} \Biggl( \alpha = 2 + 4  \myrv{y} + 10  \myrv{w}, ~ \beta = 22 - 4  \myrv{y} - 10  \myrv{w} \Biggr).
\end{split}
\end{equation}
This mechanism is also depicted visually in Figure~\ref{fig:diabetes}. Notice how, once the value of the body-mass index, $\myrv{x}_{c}$, has been established, the probability of diabetes, $\myrv{y}$, remains constant regardless of how the patient's salt intake, $\myrv{w}$, might change. This makes body-mass index, $\myrv{x}_{c}$, a consistent predictor across the environments defined by the values of the patient's salt intake, $\myrv{w}$. This is not the case for the frequency of thirst, $\myrv{x}_{s}$. For a given level of thirst frequency, $\myrv{x}_{s}$, the probability of diabetes, $\myrv{y}$, depends on the patient's salt intake, $\myrv{w}$.

The state-of-the-art approach for identifying causal features is by learning \textit{invariant predictors}. This means training the predictor with both $\myrv{x}_{c}$ and $\myrv{x}_{s}$ such that it predicts with the same accuracy across every environment. Under a certain set of assumptions, invariant predictors are guaranteed to learn to only use information from causal features, $\myrv{x}_{c}$, and to discard information from symptoms, $\myrv{x}_{s}$. I review the literature on invariant predictors in Section~\ref{sec:invariant-predictors}.

\textit{If} we knew the patient's salt intake, $W$, then the frequency of thirst, $X_2$, would become usable as a predictor. $X_2$ contains useful information about the response, $\myrv{y}$, that can potentially reduce the variance of our predictor. We can see in Figure~\ref{fig:x2-shift} that, when all other variables are kept constant, the frequency of thirst is a good predictor of diabetes.

\begin{figure}[t]
    \begin{center}
        \begin{subfigure}[t]{0.325\textwidth}
            \centerline{\includegraphics[width=\columnwidth]{figures/X1_shift.pdf}}
            \caption{}
            % \caption{Probability of diabetes, $\myrv{y}$, given body-mass, $\myrv{x}_{c}$, stays constant regardless of salt intake, $\myrv{w}$.}
            \label{fig:x1-shift}
        \end{subfigure}
        \hfill
        \begin{subfigure}[t]{0.325\textwidth}
            \centerline{\includegraphics[width=\columnwidth]{figures/X2_shift.pdf}}
            \caption{}
            % \caption{Probability of diabetes, $\myrv{y}$, given thirst frequency, $\myrv{x}_{s}$, changes with salt intake, $\myrv{w}$.}
            \label{fig:x2-shift}
        \end{subfigure}
        \hfill
        \begin{subfigure}[t]{0.325\textwidth}
            \centerline{\includegraphics[width=\columnwidth]{figures/X2_pred.pdf}}
            \caption{}
            % \caption{Probability of diabetes, $\myrv{y}$, changes with thirst frequency, $\myrv{x}_{s}$, even when salt intake, $\myrv{w}$, stays constant.}
            \label{fig:x2-pred}
        \end{subfigure}
        \caption{Causes are more consistent predictors than symptoms, although symptoms still contain important information. In (a), we can see that the probability of diabetes, $\myrv{y}$, conditioned on a cause of diabetes (body-mass index, $\myrv{x}_{c}$) remains constant across environments (i.e., regardless of whether salt intake, $\myrv{w}$, is 1 or 0). However, (b) shows that the probability of diabetes, $\myrv{y}$, conditioned on a symptom of diabetes (thirst frequency, $\myrv{x}_{s}$) changes across environments. Nevertheless, the probability of diabetes, $\myrv{y}$, changes with thirst frequency, $\myrv{x}_{s}$, within each environment (i.e., when the salt intake, $\myrv{w}$, is constant), as shown in (c).}
        \label{fig:x2}
    \end{center}
    \vskip -0.2in
\end{figure}

Crucially, we \textit{can} infer the patient's salt intake, $W$, \textit{if} we have multiple samples of the covariates, $\myrv{x}_{c}$ and $\myrv{x}_{s}$, from the same environment (i.e., from the same value of $\myrv{w}$). We can see in Figure~\ref{fig:diabetes} that multiple samples from $\myrv{x}_{c}$ would allow us to infer the distribution of $\myrv{x}_{c}$, therefore indirectly revealing the value of $\myrv{w}$.

This procedure is only possible within an evaluation paradigm that allows at test-time for multiple covariates from the same environment. In the paradigm of singleton prediction, the best we can hope for is the causal predictor. In the rest of the chapter, we will explore the two questions that arise naturally: (1) How general is this framework? (2) How would a contextual predictor be implemented in practice? At the end, I will also show an interesting link between domain adaptation and disentanglement (defined in Chapter 3).

\section{Background: Singleton predictions}

I will start by formally defining the problem of domain adaptation. Assume a training dataset comprising $n$ groups (i.e., environments), each with a different number of covariate-response pairs. Let $\myvect{x}^i_k \in \mathbb{R}^{\mydim{x}}$ be the $k$-th covariate in group $i$, with $i \in \{1,\dots,N\}$ and $k \in \{1,\dots,m_i\}$, where $\mycnt{m}_i$ is the size of group $i$. Let $\myvect{y}^i_{k} \in \mathbb{R}^{\mydim{y}}$ be the $k$-th response in group $i$. We assume that the data belonging to group $i$, the pairs $\left\{ \left( \myvect{x}^i_{k}, \myvect{y}^i_{k} \right) \right\}_{k=1}^{m_i}$, were generated by repeatedly sampling the joint distribution over two random variables, $\myrv{x}^i$ and $\myrv{y}^i$, modelling the covariate-response pairs, 
\[
    \left( \myrv{x}^i_{k}, \myrv{y}^i_{k} \right) \overset{\mathrm{i.i.d.}}{\sim} \myrv{x}^i, \myrv{y}^i, \quad \forall k \in \{1,\dots,\mycnt{m}_i\}. 
\]

The underlying joint distribution over $\myrv{x}$ and $\myrv{y}$ is assumed to change across groups. Let $\myfunc{p}^*$ represent the family of density functions corresponding to the underlying data-generating process. We allow for the general case where there may be $i \neq j \in \{1,\dots,n\}$, such that $\myfunc{p}^{*}_{\myrv{x}^i, \myrv{y}^i} \not\equiv \myfunc{p}^{*}_{\myrv{x}^j, \myrv{y}^j}$. These changes are called \textit{distributional shifts} \citep{quinonero-candelaDatasetShiftMachine2008}. 

The conventional goal of unsupervised domain adaptation is to learn a parametric density of the responses conditioned on the covariates, $\myfunc{p}^{\theta}_{\myrv{y} | \myrv{x}}$, called the \textit{predictive conditional}, that can model the data across the various groups. Let $\myfunc{p}^{\theta}_{\myrv{y} | \myrv{x}}$ be the predictive conditional parametrised by $\theta$. One possible way to quantify how well the predictor models data across multiple groups is to measure the data log-likelihood for the worst performing group,
\[
    \underset{i \in \{1,\dots,\mycnt{n}\}}{\mathrm{min}} \frac{1}{\mycnt{m}_i} \sum_{k=1}^{\mycnt{m}_i} \log \myfunc{p}^{\theta}_{\myrv{y} | \myrv{x}} \left( \myvect{y}^i_{k} | \myvect{x}^i_{k} \right).
\]

The ultimate promise of unsupervised domain adaptation lies in its ability to deliver predictors with high real-world performance, capable of accurately handling diverse and previously unseen data distributions. We assume that there is another underlying joint distribution for the in-the-wild, or held-out, data, $\myfunc{p}^{*}_{\myrv{x}', \myrv{y}'}$. Our goal is to learn a parametric distribution of the responses conditioned on the covariates, $\myfunc{p}^{\theta}_{\myrv{y} | \myrv{x}}$, that is as close as possible to the held-out distribution, $\myfunc{p}^{*}_{\myrv{y}' | \myrv{x}'}$, by using the information that can be extracted from the training data, $\left\{\left(\myvect{x}^i_{k},\myvect{y}^i_{k}\right)\right\}_{i=1, k=1}^{\mycnt{n}, \mycnt{m}_i}$.

This problem setup applies to many real-world prediction tasks in science and engineering. One example is clinical risk modelling \citep{sharmaClinicalRiskPrediction2021}. Suppose we want to create a model that assesses a patient's risk of developing cardiovascular disease, $\myrv{y}$, based on the results of an LDL cholesterol test, $\myrv{x}$. We train this model on patient case histories, where each covariate-response pair, $(\myrv{x}, \myrv{y})$ corresponds to one patient. Typically, datasets of medical records are compiled with data sourced from different environments: different hospitals, countries, years, etc. Each of these environments corresponds to a different intervention on the joint distribution, $\myfunc{p}_{\myrv{x}, \myrv{y}}$.

\subsection{Catastrophic forgetting}

One of the most surprising and problematic phenomena typical of domain adaptation is \textit{catastrophic forgetting}. The naive hope of domain adaptation is that, after we train a predictor, $\myfunc{p}^{\theta}_{\myrv{y} | \myrv{x}}$, on group $i$, training on group $j$ will improve the predictor's performance on group $i$. However, in practice, we discover that our predictor will perform poorly on group $i$, as if training on group $j$ has made the predictor ``forget'' what it had learned on group $i$. There is a large body of literature proposing improvements to neural network architectures in the hope of overcoming catastrophic forgetting in many domain areas, including medical imaging~\citep{perkoniggDynamicMemoryAlleviate2021}, machine translation~\citep{thompsonOvercomingCatastrophicForgetting2019}, or reading comprehension~\citep{xuForgetMeNot2020}.

Catastrophic forgetting is an inevitable consequence of certain kinds of distributional shifts. \citet{ben2010theory} explain that there are two kinds of shifts in the underlying joint distribution between the covariates and the responses, $\myfunc{p}^{*}_{\myrv{y} | \myrv{x}}$, with different implications for catastrophic forgetting. One is a change in the marginal distribution of the covariates (i.e., $\myfunc{p}^{*}_{\myrv{x}^i} \not\equiv \myfunc{p}^{*}_{\myrv{x}^j}$, where $i$ and $j$ are different groups) and another is a change in the predictive conditional (i.e., $\myfunc{p}^{*}_{\myrv{y}^i | \myrv{x}^i} \not\equiv \myfunc{p}^{*}_{\myrv{y}^j | \myrv{x}^j}$, where $i$ and $j$ are different groups). The former's effect on catastrophic forgetting can be mitigated by improving the neural network architecture~\citep{liuNeedLanguageDescribing2023}. The latter case, which is called conditional shift~\citep{liuNeedLanguageDescribing2023}, is more challenging. Under conditional shift, the same covariate, $\myvect{x}$, corresponds to different distributions of responses, $\myrv{y}$, depending on which environment we are in. This means that catastrophic forgetting is \textit{expected} when we train the same predictor with maximum likelihood on different environments. We must change the learning objective of the predictor in order to overcome this challenge.

\subsection{Causal predictors}
\label{sec:invariant-predictors}

Causal predictors are the gold-standard for domain adaptation under the singleton prediction paradigm. The main challenge involved in learning causal predictors is identifying which elements of the covariate, $\myrv{x}$, are causes of the response, and which are symptoms. There exists a broad literature on algorithms for discovering the causal components \citep{glymourReviewCausalDiscovery2019}.

The current state-of-the-art approach for learning causal predictors is the paradigm of \textit{invariant predictors}. A predictor is said to be invariant when it achieves the same accuracy on each of the environments in the training set. Formally, let $\myfunc{r}_i: \mathbb{R}^{\mydim{x} + \mydim{\theta}} \to \mathbb{R}$ be a function evaluating the error achieved by the predictive conditional, $\myfunc{p}^{\theta}_{\myrv{y} | \myrv{x}}$, on environment $i$. Let $\myfunc{r}_i$ be defined as:
\[
    \myfunc{r}_i \left( \myvect{x}, \theta \right) = \kl[\bigg]{\myfunc{p}^{*}_{\myrv{y}^i | \myrv{x}^i = \myvect{x}}}{\myfunc{p}^{\theta}_{\myrv{y} | \myrv{x} = \myvect{x}}}.
\]
We say that $\myfunc{p}^{\theta}_{\myrv{y} | \myrv{x}}$ is an invariant predictor if and only if the random variable $\myrv{r}_i = \myfunc{r}_i \left( \myrv{x}^i, \theta \right)$ is identically distributed across every environment: $\forall i \neq j \in \{1,\dots,\mycnt{n}\}, ~ \myrv{R}_i \sim \myrv{R}_j$. 

This definition of invariant predictors is a generalisation of the more common formulation, $\myfunc{r}_i \left( \myvect{x}, \theta \right) = \frac{1}{\mycnt{n}} \sum_{i=1}^{\mycnt{n}} \sum_{k=1}^{\mycnt{m}_i}  \left(\myvect{y}^i_k - \tilde{\myvect{y}}^i_k \right)^2$, where $\tilde{\myvect{y}}^i_k$ is the output of the predictor. I use the generalisation because, in the diabetes exaple, the response, $\myrv{y}$, is categorical rather than continuous.

Invariant predictors are attractive because, under a certain set of assumptions, they are guaranteed to learn \textit{causal features}~\citep{pearl2009causality, scholkopfCausalAnticausalLearning2012}. Causal features are the components of the covariate, $\myrv{x}$, that are direct causal parents of the response, $\myrv{y}$. A predictor trained on only causal features is guaranteed to perform with the same level of accuracy on any environment, as long as the dataset shift does not include interventions on the response, $\myrv{y}$.

There exist several algorithms for learning invariant predictors, each with its own set of specific sets of assumptions. For example, \citet{petersCausalInferenceUsing2016} assume that the relationship between the covariates, $\myrv{x}$, and responses, $\myrv{y}$, is linear. \citet{magliacane2018domain} assume that the change in the true predictive distribution, $\myfunc{p}^{*}_{\myrv{y} | \myrv{x}}$, across environments is determined by a random variable, $\myrv{u}$, that is observed. \citet{arjovskyInvariantRiskMinimization2019} assume a linear optimisation problem. 

\section{My approach: Contextual predictions}

I claim that there is another way of solving the problem of catastrophic forgetting under conditional shift: \textit{contextual prediction}. The current domain adaptation paradigm is that, at test time, we can only make predictions based on one covariate at a time. In that scenario, that invariant predictor is our best option, because the relationship between the response, $\myrv{y}$, and the components of the covariate that are symptoms of the response, $\myrv{x}_{2}$, is contingent on the environment, whereas the relationship between $\myrv{y}$ and the dimensions that are causes of the response, $\myrv{x}_{c}$, is stable across environments. What if at test time we could make use of multiple covariates from the same environment to infer something about this relationship? 

Formally, contextual prediction is an evaluation paradigm that allows the predictor to model the joint probability of responses, $\myrv{y}$, conditional on a group of covariates, $\myrv{x}$, belonging to the same environment. Let $\myfunc{p}^\theta_{\mygroup{\myrv{y}} | \mygroup{\myrv{x}}}$ be the predictive conditional density under this paradigm, where $\mygroup{\myrv{x}}$ and $\mygroup{\myrv{y}}$ are the random variables modelling a group of covariate-response pairs coming from the same environment, such as $\left\{ \left( x^i_k, y^i_k \right) \right\}_{k=1}^{\mycnt{m}_i}$.

Contextual predictors are necessary in situations where the causal predictor is not sufficient. One assumption that is necessary for invariant predictors to be a reasonable solution for domain adaptation is that the causal features of the covariate, $\myrv{x}_{c}$, must be expressive enough to predict the response, $\myrv{y}$, with any kind of accuracy. However, this is not always the case. For example, in the case of \textit{anti-causal learning}, all the components of the covariate, $\myrv{x}$, are assumed to be causal effects of the response, $\myrv{y}$, rendering invariant predictors useless~\citep{scholkopfCausalAnticausalLearning2012}.

\section{Group-instance predictor}

In this section, I will introduce a concrete example of a predictor exploiting the paradigm of contextual prediction which outperforms causal predictors in the diabetes case study introduced in Section~\ref{sec:example-diabetes}. This method is based on the group-instance model defined in the Background section. Using the group-instance model for conextual predictions is a simple idea, but it serves to show empirically the great potential of contextual predictions as an evaluation paradigm.

We first assume that we have prior knowledge of which components of the covariates, $\myrv{x}$, are causes, $\myrv{x}_{c}$, and which are symptoms, $\myrv{x}_{s}$. This assumption is fair because, for the purpose of our evaluation, the causal predictor used for comparison will be trained only on the causal components, $\myrv{x}_c$ (see Section~\ref{sec:domain-evaluation}).

We train the group-instance model to autoencode each group $i$ of symptoms $\left\{\myvect{x}^i_{s, k}\right\}_{k=1}^{\mycnt{m}_i}$ from the training set. Recall that, for each group, the group-instance encoder will infer one group-level latent vector $\myvect{w}^i \in \mathbb{R}^{\mydim{w}}$ and a group of instance-level latent vectors $\left\{ \myvect{z}^i_{k} \right\}_{k=1}^{\mycnt{m}_i}$ with $\myvect{z}^i_{k} \in \mathbb{R}^{\mydim{z}}$. Then, we train our response predictor, $\myfunc{p}^{\theta}_{\myrv{y} | \myrv{z}, \myrv{x}_c}$, the conditional density of the response, $\myrv{y}$, given the instance-level variable (inferred with the value of the corresponding symptom), $\myrv{z}$, and given the corresponding cause, $\myrv{x}_c$. This method, in effect, enhances the causal predictor by concatenating a ``de-environmentalised'' symptom to the input.

For testing, recall that the paradigm of contextual predictions means making a joint prediction for an entire group of covariates, $\left\{ \myvect{x} \right\}_{k=1}^{\mycnt{m}'}$. We use the group-instance model to infer the group-level, $\myvect{w}'$, and instance-level, $\left\{ \myvect{z}'_{k} \right\}_{k=1}^{\mycnt{m}'}$, vectors based on the held-out symptoms, $\left\{ \myvect{x}'_{k} \right\}_{k=1}^{\mycnt{m}'}$. Then, for each covariate, $k$, we obtain the distribution of the response, $\myrv{y}'_k$, from the response model, $\myfunc{p}^{\theta}_{\myrv{y} | \myrv{z} = \myvect{z}'_k, \myrv{X} = \myrv{x}'_{c}}$.

I claim that there is another way of solving the problem of catastrophic forgetting under conditional shift: \textit{contextual prediction}. The current domain adaptation paradigm is that, at test time, we can only make predictions based on one covariate at a time. In that scenario, that invariant predictor is our best option, because the relationship between the response, $\myrv{y}$, and the components of the covariate that are symptoms of the response, $\myrv{x}_{s}$, is contingent on the environment, whereas the relationship between $\myrv{y}$ and the dimensions that are causes of the response, $\myrv{x}_{c}$, is stable across environments. What if at test time we could make use of multiple covariates from the same environment to infer something about this relationship?

Formally, contextual prediction is an evaluation paradigm that allows the predictor to model the joint probability of responses, $\myrv{y}$, conditional on a group of covariates, $\myrv{x}$, belonging to the same environment. Let $\myfunc{p}^\theta_{\mygroup{\myrv{y}} | \mygroup{\myrv{x}}}$ be the predictive conditional density under this paradigm, where $\mygroup{\myrv{x}}$ and $\mygroup{\myrv{y}}$ are the random variables modelling a group of covariate-response pairs coming from the same environment, such as $\left\{ \left( x^i_k, y^i_k \right) \right\}_{k=1}^{\mycnt{m}_i}$.

Contextual predictors are necessary in situations where the causal predictor is not sufficient. One assumption that is necessary for invariant predictors to be a reasonable solution for domain adaptation is that the causal features of the covariate, $\myrv{x}_{c}$, must be expressive enough to predict the response, $\myrv{y}$, with any kind of accuracy. However, this is not always the case. For example, in the case of \textit{anti-causal learning}, all the components of the covariate, $\myrv{x}$, are assumed to be causal effects of the response, $\myrv{y}$, rendering invariant predictors useless~\citep{scholkopfCausalAnticausalLearning2012}.

The following section provides a theoretical analysis comparing the potential performance of causal and contextual predictors under a simplified, linear model that captures the essential structure of the diabetes example (Figure~\ref{fig:diabetes}). This analysis formally demonstrates the conditions under which the contextual approach can yield lower prediction error.

\section{Experiment for illustration}
\label{sec:domain-evaluation}

\begin{figure}[t]
    \begin{center}
        \centerline{\includegraphics[width=\columnwidth]{figures/domain_adaptation_results.pdf}}
        \caption{Evaluation results following the paradigm of contextual prediction. Each horizontal bar in the rugplot denotes the result for one random initialisation of the dataset and model weights. Our model, Group-Instance, which predicts jointly for an entire group of covariates, outperforms causal predictors, the current state-of-the-art under the paradigm of singleton prediction. As expected, the Naive Predictor (i.e., a model trained with maximum likelihood on the pooled data from the training environments) performs well on average but generalises poorly to new environments. The Oracle (i.e., a model similar to the Naive Predictor, but that also takes as input the ground-truth environment variable, $\myrv{w}$, in addition to the covariates, $\myrv{x}_{c}$ and $\myrv{x}_{s}$) performs best on average.}
        \label{fig:diabetes-empirical-results}
    \end{center}
    \vskip -0.4in
\end{figure}

In this section, I show empirically that the group-instance predictor outperforms the causal predictor under the diabetes case study from Section~\ref{sec:example-diabetes}. The task is to predict whether a patient has diabetes, $\myrv{y}$, based on their body-mass index, $\myrv{x}_{c}$, and their frequency of thirst, $\myrv{x}_{s}$. The patients are grouped into environments based on a hidden variable, the patient's salt intake, $\myrv{w}$.

We generate a dataset according to the diabetes mechanism defined in Equation~\ref{eq:diabetes-case-study}. Each level of salt intake, $\myrv{w} \in [0, 1, 2]$, corresponds to one data environment, two of which will be used for training and the remaining one used for testing. Specifically, $w^i = i$, for each $i \in [0,1,2]$. For each environment, $i$, we use the mechanism~\ref{eq:diabetes-case-study} to randomly generate $\mycnt{m} = \num{10000}$ covariate-response tuples containing one body-mass index, $\myvect{x}^i_{c, k}$, one frequency of thirst, $\myvect{x}^i_{s, k}$, and one diabetes diagnosis, $y^i_k$, where $k \in \{1,\dots,\mycnt{m}\}$.

The evaluation metric is the accuracy of predictions on the testing environment. We train with maximum likelihood estimation on the covariate-response pairs from the two training environments, $\left\{ \left(\myvect{x}^1_{c, k}, \myvect{x}^1_{s, k}, y^1_{k} \right) \right\}_{k=1}^{\mycnt{m}}$ and $\left\{ \left( \myvect{x}^2_{c, k}, \myvect{x}^2_{s, k}, y^2_{k} \right) \right\}_{k=1}^{\mycnt{m}}$, and evaluate on the data pairs from the testing environment, $\left\{ \left( \myvect{x}^3_{c, k}, \myvect{x}^3_{s, k}, y^3_{k} \right) \right\}_{k=1}^{\mycnt{m}}$. For each data pair, $k$, we make a prediction, $\tilde{y}^i_k$. We measure the accuracy for environment $i$ using the following function:
\[
    \mathrm{acc} \left( \theta, i \right) = \frac{1}{\mycnt{m}} \sum_{k=1}^{\mycnt{m}} \left[ \mathbb{I} \left( \tilde{y}^i_k > 0.5 \right)  y^i_k  + \mathbb{I}\left( \tilde{y}^i_k < 0.5 \right)  \left(1 - y^i_k \right) \right].
\]

I compare four prediction methods in my analysis:
\begin{enumerate}
    \item The naive predictor. The predictive conditional, $\myfunc{p}^\theta_{\myrv{y} | \myrv{x}_{c}, \myrv{x}_{s}}$, takes as input the body-mass index, $\myrv{x}_{c}$, and frequency of thirst, $\myrv{x}_{s}$, of a single patient. We implement the predictor as a binomial distribution parametrised by a linear model, $\myrv{y} = \mathrm{Bin} \left( n=1, ~ p=\theta_1 + \theta_2  \myrv{x}_{c} + \theta_3  \myrv{x}_{s} \right)$. This method is the standard supervised learning approach.
    \item The causal predictor. The predictive conditional, $\myfunc{p}^\theta_{\myrv{y} | \myrv{x}_{c}}$, takes as input only the body-mass index, $\myrv{x}_{c}$, which is the causal parent of the response, $\myrv{y}$. We implement the predictor as a binomial distribution parametrised by a linear model, $\myrv{y} = \mathrm{Bin} \left( n=1, ~ p=\theta_1 + \theta_2  \myrv{x}_{c} \right)$.  This method is the idealised predictor under the singleton prediction paradigm, since it assumes that we already know which component is causal, $\myrv{x}_c$.
    \item The group-instance predictor. The predictive conditional, $\myfunc{p}^\theta_{\myrv{y} | \myrv{x}_{c}, \myrv{x}_{s}, \mygroup{\myrv{x}}_{c}, \mygroup{\myrv{x}}_{s}}$, takes as input both the patient's covariates (i.e., body-mass index, $\myrv{x}_{c}$, and frequency of thirst, $\myrv{x}_{s}$), as well as the covariates of all the other patients in the environment, $\mygroup{\myrv{x}}_{c}$ and $\mygroup{\myrv{x}}_{s}$. We implement the predictor as a binomial distribution parametrised by a linear model, $\myrv{y} = \mathrm{Bin} \left( n=1, ~ p=\theta_1 + \theta_2  \myrv{x}_{c} + \theta_3  \myrv{z} \right)$, where $\myrv{z}$ is the inferred instance-level latent code sampled from the latent posterior density of the group-instance model, $\myfunc{q}^{\phi, \psi}_{\myrv{w}, \mygroup{\myrv{z}} | \mygroup{\myrv{x}}}$. This method is my proposed implementation for the paradigm of contextual prediction.
    \item The environment-oracle predictor. The predictive conditional, $\myfunc{p}^\theta_{\myrv{y} | \myrv{x}_{c}, \myrv{x}_{s}, \myrv{w}}$, takes as input both the patient's covariates (i.e., body-mass index, $\myrv{x}_{c}$, and frequency of thirst, $\myrv{x}_{s}$) as well as the patient's salt intake, $\myrv{w}$. This method acts as a topline because it is the only one that has access to the ground-truth salt intake levels, $w^i$. We implement the predictor as a binomial distribution parametrised by a linear model, $\myrv{y} = \mathrm{Bin} \left( n=1, ~ p=\theta_1 + \theta_2  \myrv{x}_{c} + \theta_3  \myrv{x}_{s} + \theta_4  \myrv{w} \right)$. 
\end{enumerate}

These methods are simplistic, but they reflect a real use case. Medical researchers frequently use simple linear models for clinical risk prediction \citep{nobleRiskModelsScores2011}. 

The results can be seen in Figure~\ref{fig:diabetes-empirical-results}. My group-instance model outperforms causal predictors in every environment and overall, demonstrating the potential of contextual prediction to unlock new solutions for domain adaptation. As expected, the naive predictor performs vastly differently across environments. The environment-oracle predictor performs best on average, but also has quite high variance between environments.

The results of this experiment reflect the potential of contextual predictions. Under singleton predictions, it is impossible to consistently surpass the performance of the causal predictor; it offers the best worst-case solution. However, contextual predictions allow for new methods, such as the group-instance predictor, to achieve new levels of prediction accuracy.

\section{Theoretical analysis}
\label{sec:theory-comparison}

In this section, I find a minimal set of conditions for which the group-instance predictor outperforms the causal predictor in a linear setting. To do this, I prove a theorem that shows that the group-instance predictor has higher log-likelihood on the testing data than the causal predictor \textbf{if and only if} a closed-form expression in terms of $n$ (the number of environments in the training set) and $\sigma_W$ (the standard deviation of the true environment variable) is satisfied. I consider a simplified generative model that mirrors the causal structure of the diabetes example (Figure~\ref{fig:diabetes}) but uses linear relationships and Gaussian noise. This allows for tractable calculation of prediction error variances.

\begin{assumption}[Linear Gaussian Model]
\label{ass:linear-gaussian-model}
We assume the data for any given environment $W=w$ is generated according to the following linear Gaussian structural causal model (using scalar variables for simplicity, corresponding to the core structure):
\begin{enumerate}
    \item Environment Variable: $W \sim \N(0, \sigma_W^2)$, where $\sigma_W^2 > 0$. (Analogous to $\myrv{w}$, salt intake).
    \item Causal Covariate: $X_c = \alpha_c W + \epsilon_c$, where $\epsilon_c \sim \N(0, \sigma_c^2)$. (Analogous to $\myrv{x}_c$, body-mass index).
    \item Response Variable: $Y = \beta_c X_c + \epsilon_y$, where $\epsilon_y \sim \N(0, \sigma_y^2)$, $\beta_c \neq 0$, $\sigma_y^2 > 0$. (Analogous to $\myrv{y}$, diabetes risk).
    \item Symptom Covariate: $X_s = \alpha_s W + \beta_s Y + \epsilon_s$, where $\epsilon_s \sim \N(0, \sigma_s^2)$, $\alpha_s \neq 0$, $\beta_s \neq 0$, $\sigma_s^2 > 0$. (Analogous to $\myrv{x}_s$, thirst frequency).
    \item The noise terms $\epsilon_c, \epsilon_y, \epsilon_s$ are mutually independent and independent of $W$.
    \item All model parameters ($\alpha_c, \beta_c, \alpha_s, \beta_s, \sigma_c^2, \sigma_y^2, \sigma_s^2, \sigma_W^2$) are known constants.
\end{enumerate}
\end{assumption}

We compare two predictors aiming to predict $Y_{\text{test},k}$ for a random sample $k$ from a new test environment with (unknown) environment variable $W_{\text{test}} \sim \N(0, \sigma_W^2)$. We assume access to $n$ training environments, each providing infinite samples, and infinite covariate samples $\{ (X_{c,\text{test},k}, X_{s,\text{test},k}) \}_{k=1}^\infty$ from the test environment. The goal is to minimize the expected squared prediction error, which corresponds to the variance for unbiased predictors.

\subsection{Error variance as log-likelihood}

I start by proving a simplifying result: The log-likelihood of a regression model on the testing data is greater than that of another model \textbf{if and ony if} the error variance of the first model is greater than the error variance of the second model if the models are unbiased and Linear Gaussian.

\begin{lemma}
\label{lem:loglike_vs_variance}
If the following assumptions are satisfied:
\begin{enumerate}
    \item Let $M_1$ and $M_2$ be two linear models predicting the response variable $Y$ using covariates $X_c$ and $X_s$ in a test environment generated according to Assumption~\ref{ass:linear-gaussian-model}.
    \item Let the predictions be $\hat{Y}_1$ and $\hat{Y}_2$ respectively. Assume both models are unbiased, i.e., $\E[\hat{Y}_1] = \E[Y]$ and $\E[\hat{Y}_2] = \E[Y]$ in the test environment, where the expectation is taken over the joint distribution of $(W_{\text{test}}, X_c, X_s, Y)$.
    \item Let $\sigma^2_{e_1} = \Var(Y - \hat{Y}_1)$ and $\sigma^2_{e_2} = \Var(Y - \hat{Y}_2)$ be the variances of the prediction errors.
    \item Assume further that each model $M_i$ correctly specifies its own prediction error distribution as Gaussian with the true error variance, i.e., the conditional distribution of $Y$ assumed by model $M_i$ is $p_i(Y | X_c, X_s) = \N(Y | \hat{Y}_i, \sigma^2_{e_i})$.
\end{enumerate}

Then, the expected log-likelihood of Model 1 on the test data is greater than that of Model 2 if and only if the variance of the prediction error of Model 1 is less than that of Model 2:
$$ \E[\log p_1(Y | X_c, X_s)] > \E[\log p_2(Y | X_c, X_s)] \iff \sigma^2_{e_1} < \sigma^2_{e_2} $$
where the expectation is taken over the true joint distribution of $(X_c, X_s, Y)$ in the test environment.
\end{lemma}

\begin{proof}
Let the two linear models be denoted by their predictions:
\begin{align*}
    \hat{Y}_1 &= \gamma_{c1} X_c + \gamma_{s1} X_s + \delta_1 \\
    \hat{Y}_2 &= \gamma_{c2} X_c + \gamma_{s2} X_s + \delta_2
\end{align*}
for some coefficients $\gamma_{c1}, \gamma_{s1}, \gamma_{c2}, \gamma_{s2}$ and intercepts $\delta_1, \delta_2$.

First, we determine the expected values of the variables in the test environment, where $W_{\text{test}} \sim \N(0, \sigma_W^2)$.
\begin{align*}
    \E[W_{\text{test}}] &= 0 \\
    \E[X_c] &= \E[\alpha_c W_{\text{test}} + \epsilon_c] = \alpha_c \E[W_{\text{test}}] + \E[\epsilon_c] = 0 + 0 = 0 \\
    \E[Y] &= \E[\beta_c X_c + \epsilon_y] = \beta_c \E[X_c] + \E[\epsilon_y] = 0 + 0 = 0 \\
    \E[X_s] &= \E[\alpha_s W_{\text{test}} + \beta_s Y + \epsilon_s] = \alpha_s \E[W_{\text{test}}] + \beta_s \E[Y] + \E[\epsilon_s] = 0 + 0 + 0 = 0
\end{align*}
The models are assumed to be unbiased, meaning $\E[\hat{Y}_i] = \E[Y]$.
$$ \E[\hat{Y}_i] = \E[\gamma_{ci} X_c + \gamma_{si} X_s + \delta_i] = \gamma_{ci} \E[X_c] + \gamma_{si} \E[X_s] + \delta_i = 0 + 0 + \delta_i $$
Since $\E[Y] = 0$, unbiasedness requires $\delta_i = 0$ for $i=1, 2$. Thus, the predictions are $\hat{Y}_i = \gamma_{ci} X_c + \gamma_{si} X_s$.

The prediction error for model $i$ is $e_i = Y - \hat{Y}_i$. Its expectation is $\E[e_i] = \E[Y - \hat{Y}_i] = \E[Y] - \E[\hat{Y}_i] = 0 - 0 = 0$.
The variance of the prediction error is $\sigma^2_{e_i} = \Var(e_i) = \E[(e_i - \E[e_i])^2] = \E[e_i^2] = \E[(Y - \hat{Y}_i)^2]$.

According to the lemma's assumptions, model $M_i$ uses the likelihood function corresponding to a Gaussian distribution with mean $\hat{Y}_i$ and variance $\sigma^2_{e_i}$:
$$ p_i(Y | X_c, X_s) = \frac{1}{\sqrt{2\pi \sigma^2_{e_i}}} \exp\left( -\frac{(Y - \hat{Y}_i)^2}{2 \sigma^2_{e_i}} \right) $$
The log-likelihood for a single data point $(X_c, X_s, Y)$ under model $M_i$ is:
$$ \log p_i(Y | X_c, X_s) = -\frac{1}{2} \log(2\pi) - \frac{1}{2} \log(\sigma^2_{e_i}) - \frac{(Y - \hat{Y}_i)^2}{2 \sigma^2_{e_i}} $$
We are interested in the expected log-likelihood over the true distribution of test data $(X_c, X_s, Y)$.
\begin{align*}
 \E[\log p_i(Y | X_c, X_s)] &= \E \left[ -\frac{1}{2} \log(2\pi) - \frac{1}{2} \log(\sigma^2_{e_i}) - \frac{(Y - \hat{Y}_i)^2}{2 \sigma^2_{e_i}} \right] \\
 &= -\frac{1}{2} \log(2\pi) - \frac{1}{2} \log(\sigma^2_{e_i}) - \frac{\E[(Y - \hat{Y}_i)^2]}{2 \sigma^2_{e_i}} \quad \text{(by linearity of expectation)}
\end{align*}
Using the definition $\sigma^2_{e_i} = \E[(Y - \hat{Y}_i)^2]$, we substitute this into the expression:
\begin{align*}
 \E[\log p_i(Y | X_c, X_s)] &= -\frac{1}{2} \log(2\pi) - \frac{1}{2} \log(\sigma^2_{e_i}) - \frac{\sigma^2_{e_i}}{2 \sigma^2_{e_i}} \\
 &= -\frac{1}{2} \left( \log(2\pi) + \log(\sigma^2_{e_i}) + 1 \right)
\end{align*}
Now we compare the expected log-likelihoods of Model 1 and Model 2:
$$ \E[\log p_1(Y | X_c, X_s)] > \E[\log p_2(Y | X_c, X_s)] $$
$$ \iff -\frac{1}{2} \left( \log(2\pi) + \log(\sigma^2_{e_1}) + 1 \right) > -\frac{1}{2} \left( \log(2\pi) + \log(\sigma^2_{e_2}) + 1 \right) $$
Multiply both sides by -2 and reverse the inequality sign:
$$ \iff \log(2\pi) + \log(\sigma^2_{e_1}) + 1 < \log(2\pi) + \log(\sigma^2_{e_2}) + 1 $$
Subtract the constant term $(\log(2\pi) + 1)$ from both sides:
$$ \iff \log(\sigma^2_{e_1}) < \log(\sigma^2_{e_2}) $$
Since the natural logarithm function $\log(\cdot)$ is strictly monotonically increasing, and the variances $\sigma^2_{e_1}$ and $\sigma^2_{e_2}$ are positive (as $\sigma_y^2 > 0$ and we assume non-degenerate cases where prediction isn't perfect), this inequality holds if and only if:
$$ \iff \sigma^2_{e_1} < \sigma^2_{e_2} $$
This completes the proof.
\end{proof}

\subsection{Causal vs Group-Instance predictors}

\paragraph{Causal predictor.} The Causal predictor uses only the causal covariate $X_c$, implementing the optimal linear predictor $\E[Y \given X_c]$:
\[ \hat{Y}_{C,k} = \E[Y_{\text{test},k} \given X_{c,\text{test},k}] = \beta_c X_{c,\text{test},k} .\]

The prediction error is:
\begin{align*}
    \text{Error}_{C,k} &= Y_{\text{test},k} - \hat{Y}_{C,k} \\
    &= (\beta_c X_{c,\text{test},k} + \epsilon_{y,\text{test},k}) - (\beta_c X_{c,\text{test},k}) = \epsilon_{y,\text{test},k}
\end{align*}
The variance of this error, representing the performance under the singleton prediction paradigm using only causal information, is:
\begin{equation}
    \Var(\text{Error}_{C,k}) = \Var(\epsilon_{y,\text{test},k}) = \sigma_y^2
    \label{eq:causal_var_thesis}
\end{equation}

\paragraph{Contextual predictor.} The Group-Instance predictor leverages the paradigm of contextual prediction. It uses $X_{c,\text{test},k}$, $X_{s,\text{test},k}$, and an estimate $\hat{W}_{\text{test}}$ of the unknown test environment variable $W_{\text{test}}$. This estimate is formed using the context provided by the infinite samples from the test environment (e.g., via a method like the Group-Instance model described later).

We model the estimate as $\hat{W}_{\text{test}} = W_{\text{test}} + \epsilon_W$, where $\epsilon_W$ is the estimation error with variance $\Var(\epsilon_W) = \sigma_{\epsilon_W}^2$, assumed independent of $W_{\text{test}}$ and sample-specific noise.

The optimal linear predictor using $X_c, X_s,$ and $\hat{W}$ (approximating $\E[Y \given X_c, X_s, W]$) is:
\[ \hat{Y}_{Ctx,k} = \beta_c X_{c,\text{test},k} + \Delta (X_{s,\text{test},k} - \alpha_s \hat{W}_{\text{test}} - \beta_s \beta_c X_{c,\text{test},k}) \]
where $\Delta = \frac{\beta_s \sigma_y^2}{\beta_s^2 \sigma_y^2 + \sigma_s^2}$. Let $D = \beta_s^2 \sigma_y^2 + \sigma_s^2$.

The prediction error is (derivation identical to the combined proof provided previously):
\[ \text{Error}_{Ctx,k} = (1 - \Delta \beta_s) \epsilon_{y,\text{test},k} - \Delta \epsilon_{s,\text{test},k} + \Delta \alpha_s \epsilon_W \]
The variance of the contextual predictor error is:
\begin{align}
    \Var(\text{Error}_{Ctx,k}) &= \Var((1 - \Delta \beta_s) \epsilon_{y} - \Delta \epsilon_{s}) + \Var(\Delta \alpha_s \epsilon_W) \nonumber \\
    &= \frac{\sigma_s^2 \sigma_y^2}{D} + \left(\frac{\beta_s \sigma_y^2 \alpha_s}{D}\right)^2 \sigma_{\epsilon_W}^2 \label{eq:contextual_var_general_thesis}
\end{align}
The first term, $\frac{\sigma_s^2 \sigma_y^2}{D}$, represents the optimal variance achievable if $W_{\text{test}}$ were perfectly known (the Oracle predictor variance in this linear setting). The second term represents the additional variance incurred due to the imperfect estimation of $W_{\text{test}}$.

The contextual predictor outperforms the causal predictor if and only if $\Var(\text{Error}_{Ctx,k}) < \Var(\text{Error}_{C,k})$. This depends crucially on the variance of the environment estimation error, $\sigma_{\epsilon_W}^2$. We consider two scenarios for this error variance.

\subsection{Scenario 1: Ideal contextual predictor}

This scenario assumes an efficient estimation process where the error variance scales inversely with the number of training environments, $n$.

\begin{assumption}[Ideal Estimation Error]
\label{ass:simple_error_thesis}
The variance of the environment estimation error $\epsilon_W = \hat{W}_{\text{test}} - W_{\text{test}}$ is given by:
\[ \sigma_{\epsilon_W}^2 = \Var(\epsilon_W) = \frac{\sigma_W^2}{n} \]
\end{assumption}

\begin{theorem}[Contextual vs. Causal Predictors]
\label{thm:simple_thesis}
Under the Linear Gaussian Model (Assumption \ref{ass:linear-gaussian-model}) and the Simplified Estimation Error (Assumption \ref{ass:simple_error_thesis}), the contextual predictor achieves lower prediction error variance than the causal predictor if and only if:
\[ \alpha_s^2 \frac{\sigma_W^2}{n} < \beta_s^2 \sigma_y^2 + \sigma_s^2 \]
\end{theorem}

\begin{proof}
We require $\Var(\text{Error}_{Ctx,k}) < \Var(\text{Error}_{C,k})$. Using equations \eqref{eq:causal_var_thesis}, \eqref{eq:contextual_var_general_thesis}, and Assumption \ref{ass:simple_error_thesis}:
\[ \frac{\sigma_s^2 \sigma_y^2}{D} + \left(\frac{\beta_s \sigma_y^2 \alpha_s}{D}\right)^2 \frac{\sigma_W^2}{n} < \sigma_y^2 \]
Rearranging the terms:
\[ \left(\frac{\beta_s \sigma_y^2 \alpha_s}{D}\right)^2 \frac{\sigma_W^2}{n} < \sigma_y^2 - \frac{\sigma_s^2 \sigma_y^2}{D} = \sigma_y^2 \left(1 - \frac{\sigma_s^2}{D}\right) = \sigma_y^2 \left(\frac{D - \sigma_s^2}{D}\right) \]
Substitute $D - \sigma_s^2 = \beta_s^2 \sigma_y^2$:
\[ \frac{\beta_s^2 \sigma_y^4 \alpha_s^2}{D^2} \frac{\sigma_W^2}{n} < \sigma_y^2 \left(\frac{\beta_s^2 \sigma_y^2}{D}\right) = \frac{\beta_s^2 \sigma_y^4}{D} \]
Since $\beta_s \neq 0$, $\sigma_y^2 > 0$, and $D > 0$, we can divide by $\frac{\beta_s^2 \sigma_y^4}{D}$:
\[ \frac{\alpha_s^2}{D} \frac{\sigma_W^2}{n} < 1 \]
\[ \alpha_s^2 \frac{\sigma_W^2}{n} < D \]
Substituting back $D = \beta_s^2 \sigma_y^2 + \sigma_s^2$:
\[ \alpha_s^2 \frac{\sigma_W^2}{n} < \beta_s^2 \sigma_y^2 + \sigma_s^2 \]
This completes the proof.
\end{proof}

\subsection{Scenario 2: Realistic contextual predictor}

This scenario considers a more specific estimation mechanism based on a Group-Instance Variational Autoencoder (VAE) with linear components, similar to the model proposed in this thesis and analyzed by Giordano. This provides a more precise error variance for finite $n$.

\begin{assumption}[Realistic Estimation Error.]
\label{ass:vae_error_thesis}
The estimate $\hat{W}_{\text{test}}$ is obtained via a process analogous to a linear Group-Instance VAE trained on $n$ environments ($n>2$). This process effectively estimates $W_{\text{test}}$ based on an empirical variance estimate from the training environments. The resulting estimation error variance is derived from inverse moments of the Chi-squared distribution:
\[ \sigma_{\epsilon_W}^2 = \Var(\epsilon_W) = \sigma_W^2 \left( \frac{n}{n-2} - 2 \sqrt{\frac{n}{2}} \frac{\Gamma(\frac{n-1}{2})}{\Gamma(\frac{n}{2})} + 1 \right) \quad \text{for } n > 2 \]
Let $f(n) = \left( \frac{n}{n-2} - 2 \sqrt{\frac{n}{2}} \frac{\Gamma(\frac{n-1}{2})}{\Gamma(\frac{n}{2})} + 1 \right)$. Then $\sigma_{\epsilon_W}^2 = \sigma_W^2 f(n)$. ($\Gamma(\cdot)$ is the Gamma function).
\end{assumption}

\begin{theorem}[Group-Instance vs. Causal Predictors.]
\label{thm:vae_thesis}
Under the Linear Gaussian Model (Assumption \ref{ass:linear-gaussian-model}) and the VAE-Based Estimation Error (Assumption \ref{ass:vae_error_thesis}), the contextual predictor achieves lower prediction error variance than the causal predictor if and only if (for $n>2$):
\[ \alpha_s^2 \sigma_W^2 f(n) < \beta_s^2 \sigma_y^2 + \sigma_s^2 \]
where $f(n) = \left( \frac{n}{n-2} - 2 \sqrt{\frac{n}{2}} \frac{\Gamma(\frac{n-1}{2})}{\Gamma(\frac{n}{2})} + 1 \right)$.
\end{theorem}

\begin{proof}
The proof follows the same steps as the proof of Theorem \ref{thm:simple_thesis}, substituting $\sigma_{\epsilon_W}^2 = \sigma_W^2 f(n)$ from Assumption \ref{ass:vae_error_thesis} instead of $\sigma_W^2/n$.
We require $\Var(\text{Error}_{Ctx,k}) < \Var(\text{Error}_{C,k})$.
\[ \frac{\sigma_s^2 \sigma_y^2}{D} + \left(\frac{\beta_s \sigma_y^2 \alpha_s}{D}\right)^2 \sigma_W^2 f(n) < \sigma_y^2 \]
Rearranging and simplifying as in the previous proof leads directly to:
\[ \frac{\alpha_s^2}{D} \sigma_W^2 f(n) < 1 \]
\[ \alpha_s^2 \sigma_W^2 f(n) < D \]
Substituting back $D = \beta_s^2 \sigma_y^2 + \sigma_s^2$:
\[ \alpha_s^2 \sigma_W^2 f(n) < \beta_s^2 \sigma_y^2 + \sigma_s^2 \]
This completes the proof.
\end{proof}

\subsection{Interpretation}

These theorems formalize the trade-off inherent in using contextual information.
\begin{itemize}
    \item The term $D = \beta_s^2 \sigma_y^2 + \sigma_s^2$ represents the variance of the symptom $X_s$ that is *not* directly explained by the environment $W$. It quantifies the potential reduction in prediction variance achievable by using $X_s$ if $W$ were perfectly known. A larger $D$ (stronger $Y \to X_s$ link via $\beta_s$, or larger residual symptom noise $\sigma_s^2$) indicates more potential benefit from using the symptom.
    \item The term $\alpha_s^2 \sigma_{\epsilon_W}^2$ represents the cost incurred by using the symptom $X_s$ due to the need to estimate $W$. This cost is proportional to how strongly the symptom depends on the environment ($\alpha_s^2$) and the variance of the environment estimation error ($\sigma_{\epsilon_W}^2$).
    \item The contextual predictor is superior when the potential benefit $D$ outweighs the estimation cost $\alpha_s^2 \sigma_{\epsilon_W}^2$.
    \item Both theorems show that increasing the number of training environments $n$ reduces the estimation cost (since $\sigma_W^2/n \to 0$ and $\sigma_W^2 f(n) \approx \sigma_W^2/n$ for large $n$), making the contextual predictor more likely to be advantageous. Theorem \ref{thm:vae_thesis} provides a more precise threshold for finite $n$ under the VAE estimation model.
\end{itemize}
This theoretical analysis, while based on a simplified linear Gaussian model, supports the central claim of this chapter: the paradigm of contextual prediction allows for predictors that can potentially outperform purely causal predictors by leveraging environmental context, provided the estimation of this context is sufficiently accurate. The Group-Instance predictor, described next, is one practical implementation aiming to realize this potential.

\section{Discussion}

This chapter introduced the paradigm of contextual predictions for unsupervised domain adaptation as an interesting counterpoint to singleton predictions. I start by reviewing why causal predictors are the state-of-the-art in the paradigm of singleton predictions. Causal predictors learn to predict based solely on the causes of the outcomes instead of the symptoms. This is because the relationship between cause and response does not change across environments. For example, a symptom that is positively correlated with the response in one environment may be nagatively correlated in another. This is why the fact that causal predictors ``forget'' information available in one environment is not a bug, it is a feature. 

In order to surpass the performance of the causal predictor across different environments, one way to safely make use of information from symptoms is to use the paradigm of contextual predictions. Contextual predictors make a joint prediction for an entire group of covariates belonging to the same environment. I proposed one example method for contextual predictions, but many other may exist. My method is based on the group-instance model, which comprises the core of the contribution in this thesis. This result is counterintuitive in the context of the common wisdom in domain adaptation, which states that predictions based on symptoms are unreliable.

There is one exception, however, to this argument. As I show in Chapter 3, one-sample disentanglement methods (GVAE) perform on par with my contextual disentanglement method (CxVAE), when the dataset does not suffer from conditional shift. I speculate that a similar result applies to domain adaptation: When there is no conditional shift, it is possible to design a singleton predictor that outperforms the causal predictor across multiple environments. I leave it to future work to design a singleton prediction method that effectively makes use of symptom information. A promising direction would be to use a one-sample disentanglement method, such as AdaGVAE~\citep{locatelloWeaklySupervisedDisentanglementCompromises2020} to extract a de-environmentalised representation that can be then used as input for the predictor.

\biblio

\end{document}
